% !PRINTED 237
\weberpage{237}

Es ist bei dieser Transformation stillschweigend die Voraussetzung gemacht, dass die in (5) mit $m$ und $c$ bezeichneten Gr\"ossen nicht Null seien.

Ist $m = 0$, so ist in der Formel (4) $p\,u^2 + 2\,q\,u\,v + r\,v^2$ durch $a\,u + b\,v$ theilbar, es muss dann also wenigstens f\"ur einige der Wurzeln die Gleichung
\[
a\,u + b\,v = 0,
\]
d.\,h. eine Gleichung von nicht h\"oherem als dem vierten Grade bestehen.

Ist $c = 0$, also $p\,u^2 + 2\,q\,u\,v + r\,v^2$ ein Quadrat einer linearen Function, so ergiebt sich aus (7), dass, wenn
\[
y = a'\,u + b'\,v
\]
gesetzt wird
\[
S(y) = 0, \quad S(y^2) = 0, \quad S(y^3) = 0
\]
ist, dass also die \gesperrt{Bring-Jerrard'sche} Form in rationaler Weise herstellbar ist.

Wir k\"onnen schliesslich die Resultate dieser Betrachtungen dahin zusammenfassen:

\gesperrt{Die Hauptgleichung f\"unften Grades l\"asst sich durch eine Transformation, die als einzige Irrationalit\"at die Quadratwurzel aus der Discriminante enth\"alt, auf eine Normalform mit einem Parameter transformiren.}

\begin{center}
\rule{0.25\linewidth}{0.4pt}
\end{center}
